% =====================================================================
% 02_architecture.tex — architecture et choix de conception
% =====================================================================
\chapter{Architecture et choix de conception}
\label{chap:architecture}

\section{Vue d'ensemble du modèle}

Le modèle retenu articule trois ensembles, conformément aux mécanismes
relationnels-objets décrits par Oracle \cite{oracleORD}. Le premier regroupe les acteurs et
leur organisation : agences, conseillers, clients et historique des
affectations. Le deuxième représente le patrimoine bancaire : comptes courants,
comptes d'épargne, comptes à terme et opérations. Le troisième rassemble les
mécanismes transversaux : alertes, audit, relevés mensuels et vue objet des
agences.

La figure~\ref{fig:uml} donne une vue synthétique du diagramme de classes. Les
classes y apparaissent avec leurs attributs principaux, leurs méthodes
structurantes et les associations qui les relient. Le détail exhaustif des
attributs, des contraintes et des méthodes est présenté dans les tableaux des
chapitres~\ref{chap:partieA} à~\ref{chap:partieE}. Le source PlantUML complet
reste disponible dans le fichier \code{uml/banque.puml} ; il peut être rendu
par un outil externe si nécessaire.

\begin{landscape}
\thispagestyle{plain}
\begin{center}
\vspace*{0.2cm}
\resizebox{0.88\linewidth}{!}{% =====================================================================
% uml_banque.tikz — diagramme de classes UML synthétique
% À inclure dans un environnement landscape ou une figure pleine page.
% =====================================================================
\begin{tikzpicture}[
  x=1cm, y=1cm,
  class/.style={
    draw=navy, fill=blue!4, rounded corners=2pt,
    align=left, font=\scriptsize, inner sep=4pt,
    text width=3.55cm, line width=0.6pt
  },
  note/.style={
    draw=grayText, fill=grayLight, rounded corners=2pt,
    align=left, font=\scriptsize, inner sep=3pt,
    text width=3.55cm, line width=0.4pt
  },
  assoc/.style={draw=navy, line width=0.6pt},
  inher/.style={draw=navy, line width=0.6pt, -{Triangle[open,length=6pt,width=8pt]}},
  comp/.style={draw=navy, line width=0.6pt, -{Diamond[open,length=6pt,width=8pt]}},
  dep/.style={draw=grayText, line width=0.5pt, dashed, -{Stealth[length=5pt]}},
  lbl/.style={font=\tiny, fill=white, inner sep=1pt, text=grayText}
]

% ---------------------------------------------------------------------
% Individus et structures
% ---------------------------------------------------------------------
\node[class] (agence) at (2.4,13.4) {%
  \textbf{\color{navy}agence\_type}\\[1pt]
  \textcolor{navy}{\rule{3.15cm}{0.3pt}}\\[1pt]
  code : VARCHAR2(6)\\
  nom, ville\\
  ref\_responsable : REF conseiller\_type\\[1pt]
  \textit{+ responsable()}
};

\node[class] (conseiller) at (8.4,13.4) {%
  \textbf{\color{navy}conseiller\_type}\\[1pt]
  \textcolor{navy}{\rule{3.15cm}{0.3pt}}\\[1pt]
  num\_conseiller, nom, prénom\\
  spécialité, date\_embauche\\
  salaire\\[1pt]
  \textit{+ nbClients()}
};

\node[class] (client) at (15.4,13.4) {%
  \textbf{\color{navy}client\_type}\\[1pt]
  \textcolor{navy}{\rule{3.15cm}{0.3pt}}\\[1pt]
  num\_client, nom, prénom\\
  adresse : adresse\_type\\
  téléphones : tel\_vry\_type\\
  ref\_conseiller, ref\_parrain\\[1pt]
  \textit{+ nbCepargne(), soldeTotal()}\\
  \textit{+ afficheInterets()}
};

\node[class] (affectation) at (5.4,9.6) {%
  \textbf{\color{navy}affectation}\\[1pt]
  \textcolor{navy}{\rule{3.15cm}{0.3pt}}\\[1pt]
  id\_affectation\\
  ref\_conseiller, ref\_agence\\
  date\_debut, date\_fin
};

\node[class] (compte) at (10.6,9.6) {%
  \textbf{\color{navy}compte\_type}\\[1pt]
  \textit{\color{grayText}« abstrait »}\\
  \textcolor{navy}{\rule{3.15cm}{0.3pt}}\\[1pt]
  num\_compte, solde, statut\\
  ref\_client, ref\_agence\\[1pt]
  \textit{+ fraisTenue()}\\
  \textit{+ fermeCompte() static}
};

% ---------------------------------------------------------------------
% Hiérarchie des comptes
% ---------------------------------------------------------------------
\node[class] (cptCourant) at (5.0,5.0) {%
  \textbf{\color{navy}cptCourant\_type}\\[1pt]
  \textcolor{navy}{\rule{3.15cm}{0.3pt}}\\[1pt]
  découvert, nb\_op\_cb\\
  signataires : signataire\_nt\_type\\[1pt]
  \textit{+ peutRetirer(), passeMvt()}\\
  \textit{+ nbSignataire()}
};

\node[class] (cptEpargne) at (10.6,5.0) {%
  \textbf{\color{navy}cptEpargne\_type}\\[1pt]
  \textcolor{navy}{\rule{3.15cm}{0.3pt}}\\[1pt]
  taux\_interet\\
  plafond\_depot\\[1pt]
  \textit{+ calculeInterets()}\\
  \textit{+ sortieAnticipee()}
};

\node[class] (cptTerme) at (16.2,5.0) {%
  \textbf{\color{navy}cptTerme\_type}\\[1pt]
  \textcolor{navy}{\rule{3.15cm}{0.3pt}}\\[1pt]
  montant\_bloque, duree\_mois\\
  date\_echeance, taux\_garanti\\[1pt]
  \textit{+ calculeInterets()}\\
  \textit{+ sortieAnticipee()}
};

\node[class] (operation) at (5.0,0.9) {%
  \textbf{\color{navy}operation\_type}\\[1pt]
  \textcolor{navy}{\rule{3.15cm}{0.3pt}}\\[1pt]
  id\_op, date\_op, montant\\
  type\_op, commission\\
  ref\_client, ref\_compte
};

\node[class] (signataire) at (10.2,0.9) {%
  \textbf{\color{navy}signataire\_elt\_nt\_type}\\[1pt]
  \textcolor{navy}{\rule{3.15cm}{0.3pt}}\\[1pt]
  ref\_client : REF client\_type\\
  droit : CHAR(1)\\
  \textit{nested table : signataire\_nt\_type}
};

% ---------------------------------------------------------------------
% Types valeur et collections
% ---------------------------------------------------------------------
\node[class] (adresse) at (21.4,14.2) {%
  \textbf{\color{navy}adresse\_type}\\[1pt]
  \textcolor{navy}{\rule{3.15cm}{0.3pt}}\\[1pt]
  num\_rue, rue\\
  ville, cp
};

\node[class] (telephone) at (21.4,11.7) {%
  \textbf{\color{navy}tel\_vry\_type}\\[1pt]
  \textcolor{navy}{\rule{3.15cm}{0.3pt}}\\[1pt]
  VARRAY(3) OF VARCHAR2(20)\\
  \textit{domicile, portable 1, 2}
};

\node[class] (releves) at (21.4,8.6) {%
  \textbf{\color{navy}releves}\\[1pt]
  \textcolor{navy}{\rule{3.15cm}{0.3pt}}\\[1pt]
  annee : NUMBER\\
  mois\_vry : VARRAY(12)
};

\node[class] (mois) at (21.4,6.1) {%
  \textbf{\color{navy}mois\_elt\_vry\_type}\\[1pt]
  \textcolor{navy}{\rule{3.15cm}{0.3pt}}\\[1pt]
  mois, nb\_operations\\
  types : type\_stat\_nt\_type
};

\node[class] (stat) at (21.4,3.5) {%
  \textbf{\color{navy}type\_stat\_elt\_nt\_type}\\[1pt]
  \textcolor{navy}{\rule{3.15cm}{0.3pt}}\\[1pt]
  type\_op, nb\\
  montant\_total
};

% ---------------------------------------------------------------------
% Vue objet
% ---------------------------------------------------------------------
\node[class, dashed] (vue) at (17.6,9.6) {%
  \textbf{\color{navy}« vue » Agence\_VOR}\\[1pt]
  \textcolor{navy}{\rule{3.15cm}{0.3pt}}\\[1pt]
  code, nom, ville, ref\_responsable\\
  conseillers : conseiller\_nt\_type\\
  comptes : cptCourant\_nt\_type\\[1pt]
  \textit{+ mutation()}
};

% ---------------------------------------------------------------------
% Associations
% ---------------------------------------------------------------------
% Agence - conseiller (responsable)
\draw[assoc] (agence.east) -- node[lbl,above]{1} node[lbl,below]{0..1} (conseiller.west);
% Conseiller - client (suivi)
\draw[assoc] (conseiller.east) -- node[lbl,above]{1} node[lbl,below]{0..*} (client.west);
% Agence - affectation
\draw[assoc] (agence.south) -- node[lbl,left]{1} node[lbl,right]{0..*} (affectation.north);
% Conseiller - affectation
\draw[assoc] (conseiller.south) -- node[lbl,right]{1} node[lbl,left]{0..*} (affectation.north east);
% Client - compte
\draw[assoc] (client.south) -- node[lbl,right]{1} node[lbl,left]{0..*} (compte.north east);
% Agence - compte
\draw[assoc] (agence.south east) -- node[lbl,above]{1} node[lbl,below]{0..*} (compte.north west);
% Client - adresse / téléphones
\draw[comp] (client.east) -- node[lbl,above]{1} node[lbl,below]{1} (adresse.west);
\draw[comp] (client.east) -- node[lbl,above]{1} node[lbl,below]{1} (telephone.west);
% Parrainage récursif
\draw[assoc] (client.east) .. controls +(1.3,0.7) and +(1.3,-0.7) ..
  node[lbl,right]{0..1 parrain} (client.east);
% Héritage
\draw[inher] (cptCourant.north) -- (compte.south);
\draw[inher] (cptEpargne.north) -- (compte.south);
\draw[inher] (cptTerme.north) -- (compte.south);
% Compte courant - opérations
\draw[assoc] (cptCourant.south) -- node[lbl,left]{1} node[lbl,right]{0..*} (operation.north);
% Compte courant - signataires (nested table)
\draw[comp] (cptCourant.south east) -- node[lbl,above]{1} node[lbl,below]{0..*} (signataire.north west);
% Relevés multi-niveaux
\draw[comp] (releves.south) -- (mois.north);
\draw[comp] (mois.south) -- (stat.north);
% Vue objet
\draw[dep] (vue.west) -- node[lbl,above]{vue} (compte.east);

% ---------------------------------------------------------------------
% Légende
% ---------------------------------------------------------------------
\node[note, text width=7.5cm] at (13.2,2.0) {%
  \textbf{Légende.} Trait plein : association ; losange : composition ;
  triangle ouvert : héritage ; pointillés : dépendance de vue.
  Le détail des attributs et des méthodes est donné dans les tableaux
  du chapitre 2 et dans le source PlantUML \texttt{uml/banque.puml}.
};

\end{tikzpicture}
}
\captionof{figure}{Diagramme de classes UML synthétique du projet bancaire.}
\label{fig:uml}
\end{center}
\end{landscape}

La lecture du diagramme met en évidence quatre décisions structurantes :
l'utilisation de tables objet pour les entités fortement typées, le recours à
une table relationnelle pour l'historique des affectations, l'emploi de
collections pour les données multivaluées et la présence d'une vue objet pour
agréger les informations d'une agence. Ces décisions sont justifiées dans les
sections suivantes.

\section{Tables objet et tables relationnelles : le partage}

Le projet n'oppose pas le relationnel et l'objet : il combine les deux selon
la nature des données. Les entités qui portent des références, des attributs
structurés ou des méthodes sont implantées sous forme de tables objet. Les
journaux, les historiques et les agrégats, qui se prêtent mieux à un schéma
relationnel classique, restent des tables relationnelles.

\begin{table}[htbp]
  \centering
  \small
  \caption{Partage entre tables objet et tables relationnelles.}
  \label{tab:partage-objet-relationnel}
  \begin{tabularx}{\textwidth}{@{}L{3.6cm}L{4.8cm}X@{}}
    \toprule
    \entete{Élément} & \entete{Implémentation} & \entete{Justification} \\
    \midrule
    \code{conseiller}, \code{agence}, \code{client} &
    Tables objet (\code{CREATE TABLE ... OF ...}) &
    Entités avec des \code{REF}, des méthodes et des attributs structurés. \\
    \code{cpt\_courant}, \code{cpt\_epargne}, \code{cpt\_terme} &
    Tables objet, une par sous-type &
    Le sujet interdit le compte générique ; chaque sous-type possède ses
    attributs propres. \\
    \code{operation} &
    Table objet &
    Elle référence un compte courant par \code{REF} et porte la logique de
    commission. \\
    \code{affectation} &
    Table relationnelle &
    Association datée, avec plusieurs lignes possibles par conseiller ;
    l'historique prime sur l'identité objet. \\
    \code{alerte}, \code{"AUDIT"} &
    Tables relationnelles &
    Journaux et alertes sans besoin de \code{REF} ni de méthodes. \\
    \code{releves} &
    Table relationnelle avec colonne \code{VARRAY} &
    Agrégat annuel multi-niveaux, qui n'est pas une entité objet autonome. \\
    \code{signataires} &
    \emph{Nested table} dans \code{cptCourant\_type} &
    Collection d'objets \code{signataire\_elt\_nt\_type}, stockée dans
    \code{signataire\_nt\_st}. \\
    \bottomrule
  \end{tabularx}
\end{table}

\begin{remarquebox}[Pourquoi ne pas tout faire en relationnel ?]
Le sujet impose les types objet, l'héritage, les \code{REF}, les collections,
les méthodes, le type récursif et les vues objet. Un modèle purement
relationnel ne permettrait pas de démontrer ces mécanismes. Inversement, tout
modéliser en objet serait artificiel : un historique d'affectations ou un
journal d'audit se décrit naturellement par des lignes et des colonnes. Le
partage retenu cherche donc le meilleur outil pour chaque besoin.
\end{remarquebox}

\section{Héritage des comptes}

La super-classe \code{compte\_type} est déclarée \code{NOT FINAL} et
\code{NOT INSTANTIABLE}. La première mention autorise l'héritage ; la seconde
interdit la création d'un compte générique. Cette double contrainte traduit
directement une règle métier : une banque ne gère jamais un « compte » abstrait,
mais toujours un compte courant, un compte d'épargne ou un compte à terme.

\begin{lstlisting}[caption={Déclaration de la super-classe \code{compte\_type}.},label={lst:compte-type}]
CREATE TYPE compte_type AS OBJECT (
  num_compte     VARCHAR2(10),
  solde          NUMBER(12,2),
  date_ouverture DATE,
  statut         VARCHAR2(10),
  ref_client     REF client_type,
  ref_agence     REF agence_type,
  MEMBER FUNCTION fraisTenue RETURN NUMBER
) NOT FINAL NOT INSTANTIABLE;
\end{lstlisting}

Les sous-types sont \code{cptCourant\_type}, \code{cptEpargne\_type} et
\code{cptTerme\_type}. La méthode \code{fraisTenue} est déclarée dès la
création de la super-classe, puis redéfinie dans les sous-types avec le mot-clé
\code{OVERRIDING}. Ce détail n'est pas cosmétique : Oracle ne permet pas
d'ajouter un \emph{override} par un simple \code{ALTER TYPE ... ADD MEMBER}.
Une tentative de ce type produit les erreurs \code{ORA-22324} ou
\code{PLS-00636}. La méthode est donc prévue dès la conception.

\begin{attentionbox}[Point technique]
La méthode \code{fraisTenue} est polymorphe : le même appel
\code{c.fraisTenue()} renvoie 1\,200~DA pour un compte courant et 0~DA pour un
compte d'épargne ou un compte à terme. Le comportement dépend du type réel de
l'instance, et non d'un test conditionnel écrit dans le code appelant.
\end{attentionbox}

\section{Références et intégrité}

Un \code{REF} est un pointeur vers un objet. Oracle vérifie les clés étrangères
déclarées sur des colonnes ou des attributs de type \code{REF}. Cette
vérification est essentielle : elle garantit qu'un compte ne référence pas un
client inexistant et qu'une opération ne pointe pas vers un compte supprimé.

Le modèle comporte deux références structurantes. La première est
\code{client\_type.\allowbreak ref\_parrain}, de type
\code{REF client\_type} : elle définit
un parrainage récursif, un client pouvant parrainer d'autres clients. La seconde
est \code{agence\_type.\allowbreak ref\_responsable}, de type
\code{REF conseiller\_type} :
chaque agence possède un responsable identifié.

L'historique des affectations est porté par la table \code{affectation}, qui
relie un conseiller et une agence avec une date de début et une date de fin.
Cette table permet plusieurs lignes datées pour un même conseiller et une ligne
ouverte lorsque \code{date\_fin} est nulle. Un attribut unique
\code{ref\_agence} dans \code{conseiller\_type} aurait écrasé l'historique à
chaque mutation.

\begin{clebox}[Règle de conception]
Toute relation qui doit conserver une chronologie est représentée par une table
d'association datée. Toute relation qui exprime une appartenance stable peut
être portée par un \code{REF}. Cette distinction explique le choix de
\code{affectation} pour les mutations et de \code{ref\_conseiller} pour le
conseiller de rattachement d'un client.
\end{clebox}

\section{Collections}

Le projet utilise plusieurs formes de collections, chacune adaptée à un usage
précis. Les téléphones d'un client sont bornés à trois : un \code{VARRAY(3)}
impose la limite et permet l'accès par indice. Les signataires d'un compte
courant ne sont pas bornés : une \emph{nested table} est plus appropriée. Les
statistiques mensuelles combinent un \code{VARRAY} de douze mois et une
\emph{nested table} de types d'opérations, ce qui produit une structure
multi-niveaux.

\begin{table}[htbp]
  \centering
  \small
  \caption{Collections utilisées dans le projet.}
  \label{tab:collections}
  \footnotesize
  \begin{tabularx}{\textwidth}{@{}L{2.7cm}L{5.8cm}L{3.4cm}X@{}}
    \toprule
    \entete{Collection} & \entete{Type} & \entete{Stockage} & \entete{Utilisation} \\
    \midrule
    Téléphones &
    \code{tel\_vry\_type VARRAY(3) OF VARCHAR2(20)} &
    Inline dans \code{client} &
    Jusqu'à trois numéros, éléments éventuellement \code{NULL}. \\
    Signataires &
    \code{signataire\_nt\_type TABLE OF
    signataire\_elt\allowbreak\_nt\_type} &
    \code{NESTED TABLE ... STORE AS signataire\_nt\_st} &
    Droits \code{X}, \code{R} et \code{D} sur un compte courant. \\
    Statistiques &
    \code{type\_stat\_nt\_type TABLE OF
    type\_stat\_elt\allowbreak\_nt\_type} &
    Dans \code{mois\_elt\_vry\_type} &
    Nombre et montant par type d'opération. \\
    Mois &
    \code{mois\_vry\_type VARRAY(12) OF mois\_elt\_vry\_type} &
    Dans \code{releves} &
    Douze mois d'un relevé annuel. \\
    Conseillers d'une agence &
    \code{conseiller\_nt\_type TABLE OF conseiller\_type} &
    Vue objet \code{Agence\_VOR} &
    Collection calculée par \code{CAST(\allowbreak MULTISET(...))}. \\
    Comptes d'une agence &
    \code{cptCourant\_nt\_type TABLE OF cptCourant\_type} &
    Vue objet \code{Agence\_VOR} &
    Collection calculée par \code{CAST(\allowbreak MULTISET(...))}. \\
    \bottomrule
  \end{tabularx}
\end{table}

\begin{remarquebox}[VARRAY ou \emph{nested table} ?]
Le choix dépend de la borne et de l'usage. Un \code{VARRAY} est adapté lorsque
la taille maximale est connue et faible, comme les trois téléphones. Une
\emph{nested table} est préférable lorsque le cardinal est variable, comme les
signataires. Le \code{VARRAY} autorise l'accès par indice ; la
\emph{nested table} se manipule plus naturellement avec \code{TABLE(...)} et
les opérations d'ensemble.
\end{remarquebox}

\section{Méthodes et polymorphisme}

Les méthodes sont réparties entre les types selon leur responsabilité. La
classe \code{client\_type} porte \code{nbCepargne}, \code{soldeTotal},
\code{afficheInterets} et \code{supprimeTel}. La classe
\code{cptCourant\_type} porte deux surcharges de \code{nbSignataire}, ainsi que
\code{estTitulaire}, \code{estSignataire}, \code{supprimeSign},
\code{peutRetirer} et \code{passeMvt}. Les classes \code{cptEpargne\_type} et
\code{cptTerme\_type} portent \code{calculeInterets} et
\code{sortieAnticipee}. La procédure \code{fermeCompte} est statique dans
\code{compte\_type} ; la méthode \code{mutation} appartient au type de la vue
objet \code{agence\_vor\_type}.

\begin{table}[htbp]
  \centering
  \small
  \caption{Types de méthodes et règles d'appel.}
  \label{tab:types-methodes}
  \begin{tabularx}{\textwidth}{@{}L{2.7cm}L{5.2cm}X@{}}
    \toprule
    \entete{Forme} & \entete{Exemple} & \entete{Caractéristique} \\
    \midrule
    Méthode membre &
    \code{c.nbCepargne()} &
    Appelée sur une instance ; elle peut lire \code{SELF}. \\
    Méthode statique &
    \code{compte\_type.\allowbreak fermeCompte('CC1')} &
    Appelée sur le type ; elle ne dépend pas d'une instance. \\
    Sous-programme catalogué &
    \code{virement(...)}, \code{drop\_client(...)} &
    Procédure ou fonction autonome, utilisable hors d'un type. \\
    Méthode de vue objet &
    \code{v.mutation(...)} &
    Méthode portée par le type de la vue, appliquée à une ligne de la vue. \\
    \bottomrule
  \end{tabularx}
\end{table}

\section{Triggers : où et pourquoi}

Six triggers assurent les contrôles et la traçabilité. Le trigger
\code{trg\_op\_avant} est un trigger \code{COMPOUND} : il vérifie le statut du
compte, les droits du signataire, le plafond journalier et le découvert. Le
trigger \code{trg\_op\_apres} met à jour le solde, le nombre d'opérations CB et
crée les alertes. Les quatre autres triggers tracent les échéances et les
modifications de taux ou de découvert.

\begin{table}[htbp]
  \centering
  \small
  \caption{Triggers du projet et responsabilités.}
  \label{tab:triggers-architecture}
  \begin{tabularx}{\textwidth}{@{}L{4.6cm}L{4.9cm}X@{}}
    \toprule
    \entete{Trigger} & \entete{Événement} & \entete{Rôle} \\
    \midrule
    \code{trg\_op\_avant} &
    \code{BEFORE INSERT} sur \code{operation}, \code{COMPOUND} &
    Vérifie le statut, les droits, le plafond journalier et le découvert. \\
    \code{trg\_op\_apres} &
    \code{AFTER INSERT} sur \code{operation} &
    Met à jour le solde et \code{nb\_op\_cb}, crée les alertes. \\
    \code{trg\_ct\_echeance} &
    \code{AFTER INSERT OR UPDATE OF date\_echeance} sur \code{cpt\_terme} &
    Crée une alerte \code{ECHEANCE\_PROCHE} ou \code{ECHEANCE\_DEPASSEE}. \\
    \code{trg\_audit\_taux} &
    \code{AFTER UPDATE OF taux\_interet} sur \code{cpt\_epargne} &
    Trace l'ancien et le nouveau taux dans \code{"AUDIT"}. \\
    \code{trg\_audit\_taux\_terme} &
    \code{AFTER UPDATE OF taux\_garanti} sur \code{cpt\_terme} &
    Trace l'ancien et le nouveau taux garanti. \\
    \code{trg\_audit\_decouvert} &
    \code{AFTER UPDATE OF decouvert} sur \code{cpt\_courant} &
    Trace l'ancien et le nouveau découvert. \\
    \bottomrule
  \end{tabularx}
\end{table}

\begin{attentionbox}[Table mutante]
Un trigger ligne par ligne qui interroge \code{operation}, la table en cours de
modification, provoque l'erreur \code{ORA-04091}. Le trigger \code{COMPOUND}
résout ce problème en séparant les phases : \code{BEFORE STATEMENT} prend un
instantané des retraits du jour, \code{BEFORE EACH ROW} accumule le montant de
la ligne courante, et la partie ligne par ligne ne relit jamais
\code{operation}. Ce mécanisme est détaillé au
chapitre~\ref{chap:partieD2}.
\end{attentionbox}

\section{Vue objet des agences}

La vue \code{Agence\_VOR} est une vue objet de type \code{agence\_vor\_type}.
Elle expose les attributs de l'agence, une collection de conseillers
actuellement affectés, une collection de comptes courants et une méthode
\code{mutation}. Elle est construite avec \code{CAST(MULTISET(...))} et
\code{MAKE\_REF(agence, a.code)}. Le recours à \code{MAKE\_REF} est possible
parce que la table \code{agence} possède un
\code{OBJECT IDENTIFIER IS PRIMARY KEY}. Ce point est développé au
chapitre~\ref{chap:partieE}.

\section{Synthèse des choix à défendre}

Le tableau~\ref{tab:choix-a-defendre} rassemble les choix qui peuvent être
interrogés en soutenance. Chaque ligne présente la décision, sa justification
et l'alternative écartée. Cette forme est volontairement concise : elle permet
de répondre rapidement, puis de développer si le jury insiste.

\begin{table}[htbp]
  \centering
  \small
  \caption{Choix de conception et alternatives rejetées.}
  \label{tab:choix-a-defendre}
  \begin{tabularx}{\textwidth}{@{}L{3.8cm}X L{4.1cm}@{}}
    \toprule
    \entete{Choix} & \entete{Justification} & \entete{Alternative rejetée} \\
    \midrule
    Une table par sous-type de compte &
    Chaque sous-type a ses attributs ; pas de colonnes nulles inutiles ; le
    sujet interdit le compte générique. &
    Une seule table \code{compte} avec un discriminant et des colonnes
    nullables. \\
    \code{client} avec OID système &
    Un \code{REF} basé sur la clé primaire ne peut pas être inséré dans une
    \emph{nested table} (\code{ORA-22979}). &
    \code{OBJECT IDENTIFIER IS PRIMARY KEY} sur \code{client}, qui bloque les
    signataires. \\
    \code{agence} avec OID basé sur la clé primaire &
    Nécessaire pour \code{MAKE\_REF(agence, a.code)} dans la vue objet. &
    OID système, qui rend \code{MAKE\_REF} impossible. \\
    \code{affectation} relationnelle &
    Historique daté, plusieurs lignes par conseiller. &
    Attribut \code{ref\_agence} unique dans \code{conseiller\_type}. \\
    \code{ref\_parrain} ajouté par \code{ALTER TYPE} &
    Le sujet demande le parrainage ; cela montre aussi l'évolution d'un type. &
    Recréer \code{client\_type} de zéro, avec perte des dépendances. \\
    \code{CHECK(ref\_client IS NOT NULL) DEFERRABLE} sur \code{cpt\_courant} &
    B5b insère un compte courant sans titulaire, puis le met à jour. &
    \code{NOT NULL} immédiat, qui rend B5b impossible. \\
    Trigger sur \code{operation}, pas sur la \emph{nested table} &
    Oracle interdit les triggers sur le stockage d'une \emph{nested table}
    (\code{ORA-22812}). &
    Trigger sur \code{signataire\_nt\_st}, impossible. \\
    \code{virement} sans \code{COMMIT} &
    L'appelant doit pouvoir enchaîner ou annuler. &
    \code{COMMIT} interne, qui empêche de tester le rollback global. \\
    \code{fermeCompte} avec \code{COMMIT} &
    C'est une opération de clôture définitive, testée en dernier. &
    Absence de \code{COMMIT}, qui ne persiste pas la suppression. \\
    \bottomrule
  \end{tabularx}
\end{table}

\section{Conclusion du chapitre}

L'architecture retenue n'est ni purement relationnelle ni purement objet. Elle
utilise l'objet là où il apporte une réelle valeur — héritage, polymorphisme,
références, collections, méthodes — et le relationnel là où il reste plus
simple et plus lisible. Cette hybridation est conforme au sujet et facilite la
démonstration : chaque mécanisme peut être montré isolément, puis intégré dans
un scénario complet.
